\documentclass[bachelor,assignment]{pmk-graduate}


\begin{document}
% Учебный год
\year{2025-2026}

% Приказ: дата, номер
%\order{02.12.2025}{222-С}

% Студент
\student{Милов Данила Константинович
\brief Д.\,К.\,Милов
\group 47
\english Milov Danila Konstantinovich
}

% Научный руководитель
\superviser{Дудаков Сергей Михайлович
\brief С.\,М.\,Дудаков
\position декан факультета ПМиК
\degree д.\,ф.-м.\,н. % Учёная степень. При отсутствии убрать строку
\rank профессор % Учёное звание. При отсутствии убрать строку
}

% Образовательная программа. См. перечень ниже
\program{ics}
% Бакалавриат:
% ics - Информатика и компьютерные науки
% swi - Инженерия программного обеспечения
% swiai - Программная инженерия в искусственном интеллекте
% aida - Искусственный интеллект и анализ данных
% mm - Математическое моделирование
% sa - Системный анализ
% icpim - Инженерное компьютерное проектирование и индустриальная математика
% acse - Прикладная информатика в экономике
% acsm - Прикладная информатика в мехатронике
% icmrs - Интеллектуальное управление в мехатронных и робототехнических системах


% Заголовок ВКР
\title{Изучение методов фрактального сжатия для различных типов информации
    \english Fractal compression of different types of information}

% Цель ВКР
\goal{Изучить и реализовать алгоритмы фрактального сжатия для различных типов информации.}

% Задачи для разработки в ВКР
\begin{tasks}
  \task{Изучение алгоритмов фрактального сжатия для изображений.
  }
  \task{Изучение алгоритмов фрактального сжатия звука.
  }
  \task{Реализация алгоритмов фрактального сжатия.
  }
  \task{Сравнение качества сжимающих алгоритмов.
  }
  \task{Адаптация алгоритмов фрактального сжатия для других типов информации.
  }
\end{tasks}

% Перечень иллюстративного материала
%\begin{images}
%\image{1}    
%\image{2}    
%\end{images}

% Консультанты
%\begin{advisers}
%\adviser{Сидоров Сидор Сидорович, проектирование схемы БД}    
%\end{advisers}

% Команда для генерации задания на ВКР
\makeassignment
\end{document}
